% main.tex - articulo unificado (tesis-bsf-co2-final)
% Compilar: xelatex main; biber main; xelatex main; xelatex main
% metodos.tex queda como referencia legacy; no se incluye.
\documentclass[11pt,a4paper]{article}

\usepackage{fontspec}
\setmainfont{TeX Gyre Termes}
\setsansfont{TeX Gyre Heros}
\setmonofont{Latin Modern Mono Light}

\usepackage[spanish,es-noquoting,es-nodecimaldot]{babel}
\usepackage{amsmath}
\usepackage[margin=2.5cm]{geometry}
\usepackage{graphicx}
\usepackage{booktabs}
\usepackage{siunitx}
\usepackage{caption}
\usepackage{lineno}
\linenumbers

\usepackage[backend=biber,style=authoryear,sorting=nyt]{biblatex}
\addbibresource{bibliografia/bibliografia.bib}

\usepackage[hidelinks]{hyperref}

\title{Modelado de la producci\'on de CO$_2$ en la bioconversi\'on con
  \textit{Hermetia illucens}: separaci\'on de fuentes larvaria y
  microbiana en c\'amaras est\'aticas}
\author{%
\begin{tabular}{cc}
John Jairo Leal G\'omez${}^{1}$ & Sorany Milena Barrientos${}^{2}$\\[2pt]
Luis Octavio Gonz\'alez${}^{1}$ & Luis Fernando Mej\'ia${}^{1}$
\end{tabular}\\[8pt]
{\small ${}^{1}$Facultad de Ingenier\'ia y Administraci\'on}\\
{\small ${}^{2}$Facultad de Ciencias Agropecuarias}\\
{\small Universidad Nacional de Colombia, Sede Palmira, Valle del Cauca, Colombia}
}
\date{\today}

\begin{document}
\maketitle

% TODO: resumen (<=300 palabras) y keywords -- Paso 6

\section{Introducci\'on}
\label{sec:intro}

\section{Materiales y m\'etodos}
\label{sec:metodos}

\subsection{Sistema experimental y protocolo de medición}
\label{sec:sistema}

Las mediciones de gas se realizaron en cámaras plásticas rígidas de geometría troncopiramidal (volumen geométrico $\approx 13.8$~L), selladas durante cada cierre con su tapa reforzada con silicona, de modo que la concentración interior de CO$_2$ solo cambia por la producción del contenido. El sensor de CO$_2$ (MH-410D), montado en una cara lateral sobre el nivel del sustrato, se verificó contra un cromatógrafo de gases con fuente certificada en el rango del experimento (0--6000~ppm). El sustrato y las larvas se alojaron en un recipiente interno con tapa de malla, bajo el supuesto de equilibrio rápido de gases con el aire de la cámara; el mismo tipo de montaje se usó en tratamientos y controles, y el volumen de aire efectivo $V_{\mathrm{aire}}$ ---el volumen geométrico menos el ocupado por el recipiente y su contenido--- entró como parámetro medido en el cálculo de tasas. Las dimensiones completas del montaje se reportan en el Material Suplementario (sec.~\ref{sup:montaje}).

Se evaluaron dos dietas, D1 y D4, cuya descripción y composición proximal se
detallan en el Cuadro~\ref{tab:dietas}.

\begin{table}[htbp]
  \centering
  \caption{Dietas evaluadas: descripción y composición proximal del análisis
  bromatológico propio (\% en base seca; EE: extracto etéreo).}
  \label{tab:dietas}
  \begin{tabular}{lp{5.4cm}ccc}
    \hline
    Dieta & Descripción & Proteína & EE & Cenizas \\
    \hline
    D1 & Alimento comercial para pollo de engorde (dieta de referencia del
         modelo de Eriksen) & 21.4 & 4.2 & 8.4 \\
    D4 & Residuos agroindustriales de frutas (naranja, plátano, entre
         otras), fermentados & 13.0 & 8.1 & 9.3 \\
    \hline
  \end{tabular}
\end{table}

Cada lote recibió dos cargas de alimento, sin retiros: 500~g por recipiente a la siembra y 250~g al día~5, tal como se recibió, sin ajuste de humedad (materia seca del análisis bromatológico: 86.3~\% en D1 y 80.3~\% en D4), con 700 larvas por recipiente; la equivalencia en materia seca de cada carga se reporta en el Material Suplementario (sec.~\ref{sup:montaje}). El experimento se dio por finalizado el día~15. Como controles se incluyeron recipientes con alimento solo (el mismo régimen de cargas, sin larvas; respiración del sustrato) y con larvas solas (mantenimiento sin alimento), medidos con el mismo protocolo; la tasa neta de cada tratamiento se obtuvo restando el control correspondiente del mismo día. La temperatura de la cámara se registró de forma continua y se mantuvo aproximadamente constante entre cierres.

En cada día de medición los cierres se hicieron en secuencia por dieta, con los dos recipientes de cada dieta medidos uno tras otro; cada cierre se prolongó hasta la saturación del sensor o, en su defecto, durante una ventana fija de 15~min con registros cada 30~s. Cada unidad se midió con tres cierres consecutivos sobre el mismo lote (réplicas técnicas de medición); el valor reportado por punto es la media de sus tres cierres y el rango su dispersión. El experimento de gas no incluyó réplicas biológicas (lotes independientes).

La tasa de cada cierre se estimó por ajuste lineal del tramo previo a la saturación; cuando la curva alcanzó el límite superior del rango del sensor, la tasa estimada subestima la tasa real y se trató como cota inferior en la validación (sec.~\ref{sec:resultados}).

En un ensayo paralelo, con las mismas condiciones de cría, se siguió el crecimiento larvario en cinco dietas (D1--D5), con tres recipientes (réplicas) por dieta, 700 larvas por recipiente y el mismo régimen de alimentación que las unidades de gas. En los días 0, 2, 4, 6, 8, 10 y 14 se tomó de cada recipiente una muestra de 10 larvas, que se pesó en conjunto de forma no destructiva y se devolvió al recipiente; cada punto es la masa fresca media por larva (peso de la muestra entre diez) y su dispersión la desviación estándar entre las tres réplicas. Las curvas de D1 y D4, las dietas del experimento de gas, entraron en el ajuste conjunto del modelo (sec.~\ref{sec:calibracion}).



\subsection{Modelo}
\label{sec:modelo}

\subsection{Calibraci\'on y protocolo de evaluaci\'on}
\label{sec:calibracion}

\section{Resultados}
\label{sec:resultados}

\section{Discusi\'on}
\label{sec:discusion}

\section{Conclusiones}
\label{sec:conclusiones}

% Material suplementario - secciones S1..S9
% Contenido de cada una se escribe cuando lleguemos a su tema.
\section*{Material suplementario}

\subsection{Balance molar de la c\'amara y conversi\'on a ppm}
\label{sup:balance}
% TODO S1

\subsection{Extensiones del n\'ucleo larvario al r\'egimen por lotes}
\label{sup:extensiones}
% TODO S2: factor de calidad, limitacion por alimento, Bmax(t), m prepupa

\subsection{M\'odulo microbiano del lecho: constante efectiva y humedad}
\label{sup:microbiano}
% TODO S3

\subsection{Justificaci\'on de las cotas de b\'usqueda de la calibraci\'on}
\label{sup:cotas}
% TODO S4

\subsection{Par\'ametros fijos y tomados de la literatura}
\label{sup:params}
% TODO S5

\subsection{Dimensiones del montaje y equivalencias de materia seca}
\label{sup:montaje}
% TODO S6: recuperar del viejo S9 (volumen de aire, cargas en MS por dieta)

\subsection{Propagaci\'on de la incertidumbre de las tasas}
\label{sup:incertidumbre}
% TODO S7

\subsection{Validaci\'on externa: digitalizaci\'on y transferencia de estados}
\label{sup:externa}
% TODO S8

\subsection{Correspondencia entre el modelo de Eriksen y la implementaci\'on}
\label{sup:correspondencia}
% TODO S9: tabla ecuacion paper <-> Excel oficial <-> modelo_bsf.py


\printbibliography

\end{document}
